
\documentclass[11pt]{article}

\usepackage[margin=1.03in]{geometry}
\usepackage{amsmath,amssymb,amsthm,mathtools}
\usepackage{microtype}
\usepackage[hidelinks]{hyperref}
\usepackage{booktabs}
\usepackage{enumitem}
\usepackage{array}

\newtheorem{theorem}{Theorem}[section]
\newtheorem{proposition}[theorem]{Proposition}
\newtheorem{lemma}[theorem]{Lemma}
\newtheorem{corollary}[theorem]{Corollary}
\newtheorem{remark}[theorem]{Remark}
\newtheorem{question}[theorem]{Question}
\newtheorem{definition}[theorem]{Definition}

\newcommand{\Area}{\operatorname{Area}}
\newcommand{\Jac}{\operatorname{Jac}}
\newcommand{\T}{\mathbb T}
\newcommand{\dd}{\,d}
\newcommand{\C}{\mathbb C}
\newcommand{\R}{\mathbb R}
\newcommand{\Z}{\mathbb Z}
\newcommand{\eps}{\varepsilon}
\newcommand{\ip}[2]{\langle #1,#2\rangle}

\title{\bfseries Fourier and Resonant Nonlinear Bounds for\\
Gromov's Filling Area Problem}
\author{}
\date{Draft, August 26, 2026}

\begin{document}
\maketitle

\begin{abstract}
Let $M$ be a compact Riemannian surface whose boundary is isometric to the
Riemannian circle of length $2\pi$, and assume that $M$ introduces no
shortcuts between boundary points.  We develop a Fourier method for lower
bounds on $\Area(M)$.

First, by taking the odd Fourier coefficients of the boundary-distance
family and applying an explicit orthogonal ``untwisting'', we obtain the
orientation-free bound
\[
        \boxed{\Area(M)\ge \frac{14\zeta(3)}{\pi}}
        =5.3567723444\ldots .
\]
The proof uses only Lipschitz distance functions, Bessel's inequality,
elementary Givens rotations, the area formula, and mod-$2$ degree.
A refinement keeps the antipodal slack and gives the exact defect estimate
\[
 \Area(M)\ge\frac{14\zeta(3)}{\pi}
 +\frac1{2\pi}\int_M\int_0^{2\pi}
 \left|\nabla\frac{
 u_\theta+u_{\theta+\pi}-\pi}{2}\right|^2
 \,d\theta\,dA .
\]

Second, in the oriented case we exploit genuinely nonlinear relations
between the Fourier modes.  If
\[
 z_n=A_n+iB_n,\qquad n=1,3,5,\ldots,
\]
we introduce the resonant deformation
\[
        G_n=z_n+\lambda\,\overline{z_1}^{\,2}z_{n+2}.
\]
On the boundary the nonlinear term has exactly the same phase as $z_n$, so
the boundary symplectic action increases linearly in $\lambda$.  At every
plane saturating the original Fourier comass, however, the first variation
of the comass vanishes by the Hardy-space autocorrelation identity
\[
        \sum_{n>0,\ n\ {\rm odd}}c_n\overline{c_{n+2}}=0.
\]
A quantitative version of this cancellation gives an explicit global
comass estimate and, for $\lambda=0.03$,
\[
        \boxed{\Area(M)>5.38982446}
\]
for every oriented Riemannian isometric filling.  Optimizing the explicit
one-parameter estimate gives approximately $5.38984897$.

We also record two barriers.  Scalar, separable $1$-Lipschitz nonlinear
reparametrizations of the distance functions cannot improve
$14\zeta(3)/\pi$, and any argument which treats the Fourier coefficient
space merely as a Euclidean Hilbert space is blocked by a holomorphic disk
of area $14\zeta(3)/\pi$.  Thus the nonlinear gain necessarily uses the
special $L^\infty$ geometry of genuine distance profiles.  Finally, we
describe an infinite hierarchy of phase-resonant perturbations whose
first-order comass cost vanishes on the entire linear saturation set.  For
the normalized one-high-leg hierarchy we prove a complementary spectral
ceiling: every certificate obtained from a global ambient comass estimate is
strictly smaller than
\[
        5.466.
\]
The ceiling is the norm of an explicit positive trace-class resonance matrix.
This shows that higher powers of the same phase-correction mechanism cannot
approach the conjectural value $2\pi$ and points instead toward profile-adapted
geometry, including a Fourier--Z\"ust hybrid and a de Sitter reformulation of
Z\"ust's kernel.

Briggs and Wells recently state that their bound
$(\sqrt3/4)\pi^2$ appears to be the first quantitative lower bound applying
to arbitrary isometric fillings.  Z\"ust records the oriented global bound
$\pi^2/2$ and proves a substantially deeper almost-calibration theorem near
the hemisphere.  To our knowledge, the estimates in this note are new.
\end{abstract}

\tableofcontents

\section{Introduction}

Let
\[
        \T=\R/2\pi\Z
\]
be the Riemannian circle of circumference $2\pi$, with intrinsic distance
\[
        d_{\T}(s,t)=\min_{k\in\Z}|s-t+2\pi k|.
\]
Let $M$ be a compact connected Riemannian surface with one boundary
component, and let
\[
        \gamma:\T\longrightarrow\partial M
\]
be an arclength parametrization.  We call $M$ an \emph{isometric filling}
of $\T$ if
\begin{equation}\label{eq:isometric}
        d_M(\gamma(s),\gamma(t))=d_{\T}(s,t)
        \qquad(s,t\in\T).
\end{equation}
The reverse inequality
$d_M(\gamma(s),\gamma(t))\le d_{\T}(s,t)$ is automatic by traveling along
the shorter boundary arc, so \eqref{eq:isometric} is precisely the
``no-shortcuts'' condition.

Gromov's filling area conjecture asserts that every \emph{orientable}
Riemannian isometric filling satisfies
\[
        \Area(M)\ge 2\pi,
\]
with equality for the unit round hemisphere \cite{Gromov1983}.  Gromov
proved the disk case.  Bangert--Croke--Ivanov--Katz proved the genus-one
case, in fact in a hyperelliptic setting
\cite{BangertCrokeIvanovKatz2005}.  The general conjecture remains open.

Briggs and Wells recently proved, without orientability,
\[
        \Area(M)\ge\frac{\sqrt3}{4}\pi^2
        =4.2736640683\ldots,
\]
and state that, as far as they are aware, this is the first quantitative
bound that applies to arbitrary isometric fillings
\cite{BriggsWells2026}.  Z\"ust records the oriented current-theoretic
projection estimate
\[
        \Area(M)\ge\frac{\pi^2}{2}
        =4.9348022005\ldots
\]
in the Riemannian setting, and proves a much finer almost-calibration
theorem near the hemisphere \cite{Zuest2025}.

The purpose of this note is to develop a direct Fourier approach.  Its
linear part is almost elementary and already gives a substantially stronger
universal estimate.  A nonlinear phase resonance then yields a further
strict improvement for the oriented conjecture.

Our first main theorem is orientation-free.

\begin{theorem}[Universal Fourier bound]\label{thm:universal}
Every compact connected Riemannian isometric filling of the Riemannian
circle of circumference $2\pi$, orientable or not, satisfies
\[
        \boxed{\Area(M)\ge\frac{14\zeta(3)}{\pi}}
        =5.3567723444\ldots .
\]
\end{theorem}

The second main theorem uses orientation.

\begin{theorem}[Explicit nonlinear improvement]\label{thm:nonlinear-main}
Every compact connected \emph{oriented} Riemannian isometric filling of
the Riemannian circle of circumference $2\pi$ satisfies
\[
        \boxed{\Area(M)>5.38982446}.
\]
More precisely, Theorem~\ref{thm:nonlinear-formula} below gives an explicit
one-parameter family of lower bounds.  Its numerical optimum is
approximately
\[
        5.38984897.
\]
\end{theorem}

Thus the successive universal/oriented constants discussed here are
\[
\begin{array}{c|c|c}
\text{bound}&\text{value}&\text{fraction of }2\pi\\ \hline
(\sqrt3/4)\pi^2
&4.2736640683\ldots&0.68017\ldots\\
\pi^2/2
&4.9348022005\ldots&0.78539\ldots\\
14\zeta(3)/\pi
&5.3567723444\ldots&0.85255\ldots\\
\text{nonlinear Fourier}
&5.38982446\ldots&0.85781\ldots\\
2\pi
&6.2831853071\ldots&1.
\end{array}
\]

\subsection*{Structure of the argument}

For each boundary point $\gamma(\theta)$ define
\[
        u_\theta(x)=d_M(x,\gamma(\theta)).
\]
The odd Fourier coefficients of the family $\theta\mapsto u_\theta(x)$
give planar maps
\[
        F_n=(A_n,B_n):M\to\R^2,\qquad n=1,3,5,\ldots.
\]
A pointwise Bessel estimate gives a common unit budget for the sum of their
Jacobian determinants.  On the boundary, the $n$th map is the $n$-fold
circle
\[
        -\frac4{\pi n^2}e^{ins}.
\]
In the orientable case, integer degree immediately converts the winding
$n$ into area.  Without orientability, a finite orthogonal mixing of the
Fourier modes turns all of these multiple circles into Jordan curves, whose
interiors are detected by mod-$2$ degree.  This proves
Theorem~\ref{thm:universal}.

The nonlinear part starts from the observation that the boundary phases
satisfy
\[
 \overline{z_1}^{\,2}z_{n+2}\sim e^{ins}.
\]
Thus the cubic monomial is resonant with the $n$th Fourier coordinate.
The boundary action therefore improves to first order.  At a plane which
saturates the original Fourier comass, the associated boundary-direction
function is an analytic or anti-analytic inner function, and all its
nonzero autocorrelations vanish.  This kills the first variation of the
comass.  Quantifying this cancellation leads to
Theorem~\ref{thm:nonlinear-main}.

The same phase algebra generates an infinite sequence of higher resonances.
Section~\ref{sec:one-high-leg-ceiling} shows, however, that this direction has
a genuine quantitative endpoint: after normalization, all one-high-leg
resonances are controlled by one positive trace-class matrix whose operator
norm is $<5.466$.  Section~\ref{sec:beyond-ceiling} then isolates the
mechanisms which can evade that spectral obstruction.

\section{Odd boundary-distance profiles}

For $\theta\in\T$, set
\begin{equation}\label{eq:u}
        u_\theta(x)=d_M(x,\gamma(\theta)).
\end{equation}
Each $u_\theta$ is $1$-Lipschitz in $x$:
\[
        |u_\theta(x)-u_\theta(y)|\le d_M(x,y).
\]
Moreover,
\[
        |u_\theta(x)-u_\varphi(y)|
        \le d_M(x,y)+d_{\T}(\theta,\varphi),
\]
so $(x,\theta)\mapsto u_\theta(x)$ is jointly Lipschitz in local
coordinates.

The odd part of the distance profile will be more convenient.

\begin{definition}\label{def:odd-slack}
Define
\begin{align}
 a_\theta(x)
 &:=\frac{u_\theta(x)-u_{\theta+\pi}(x)}{2},
 \label{eq:oddprofile}\\
 \sigma_\theta(x)
 &:=\frac{u_\theta(x)+u_{\theta+\pi}(x)-\pi}{2}.
 \label{eq:slack}
\end{align}
\end{definition}

The triangle inequality and \eqref{eq:isometric} give
\[
        \sigma_\theta(x)\ge0.
\]
Also
\[
        a_{\theta+\pi}=-a_\theta,\qquad
        \sigma_{\theta+\pi}=\sigma_\theta.
\]
On the boundary,
\begin{equation}\label{eq:boundary-odd}
        a_\theta(\gamma(s))
        =d_{\T}(\theta,s)-\frac{\pi}{2},
        \qquad
        \sigma_\theta(\gamma(s))=0.
\end{equation}

Both $\theta\mapsto a_\theta(x)$ and $x\mapsto a_\theta(x)$ are
$1$-Lipschitz.  Indeed, $a$ is the average of two differences of
$1$-Lipschitz functions.

We shall use the following standard eikonal fact in a form which is
convenient for the present argument.

\begin{lemma}[Eikonal identity for boundary distance functions]
\label{lem:eikonal-distance}
For every fixed $\theta$, the distance function
\[
        u_\theta(x)=d_M(x,\gamma(\theta))
\]
satisfies
\[
        |\nabla u_\theta(x)|=1
\]
at every interior point at which it is differentiable.  Consequently,
\[
        |\nabla_xu_\theta(x)|=1
\]
for almost every $(x,\theta)\in M^\circ\times\T$.
\end{lemma}

\begin{proof}
Since $u_\theta$ is $1$-Lipschitz, differentiability at $x$ immediately
gives
\[
        |\nabla u_\theta(x)|\le1.
\]
For the reverse inequality choose a unit-speed minimizing geodesic
\[
        \eta:[0,L]\to M,\qquad
        \eta(0)=\gamma(\theta),\quad \eta(L)=x,
\]
where $L=d_M(\gamma(\theta),x)$.  Every initial subsegment is minimizing,
so
\[
        u_\theta(\eta(t))=t,\qquad0\le t\le L.
\]
Indeed, a strict shortening to an intermediate point would concatenate
with the remaining part of $\eta$ to shorten the path to $x$.

Writing $v=\dot\eta(L)$, we therefore have
\[
 \frac{
 u_\theta(x)-u_\theta(\eta(L-h))
 }{h}=1.
\]
If $u_\theta$ is differentiable at $x$, then letting $h\downarrow0$ gives
\[
        du_\theta(x)(v)=1.
\]
Since $|v|=1$,
\[
        |\nabla u_\theta(x)|\ge1,
\]
and hence equality holds.

Rademacher's theorem gives differentiability almost everywhere for every
fixed $\theta$.  Fubini then yields the product almost-everywhere
statement.
\end{proof}

Applying Lemma~\ref{lem:eikonal-distance} to $\theta$ and
$\theta+\pi$, the parallelogram identity yields
\begin{equation}\label{eq:parallelogram-slack}
        \boxed{
        |\nabla a_\theta|^2+|\nabla\sigma_\theta|^2=1
        }
\end{equation}
for almost every $(x,\theta)$.

\begin{remark}
The map
\[
        f_x(\theta):=\frac{\pi}{2}+a_\theta(x)
\]
is anti-periodic:
$f_x(\theta+\pi)=\pi-f_x(\theta)$.
It is the canonical odd projection of the boundary-distance profile and
is closely related to the standard realization of the injective hull
$E(S^1)$ used by Z\"ust.  The slack $\sigma$ measures the part discarded
by this projection.
\end{remark}

\section{Fourier coordinates and the Jacobian budget}

For $n\ge1$ define
\begin{align}
 A_n(x)&:=\frac1\pi\int_0^{2\pi}
             a_\theta(x)\cos(n\theta)\,d\theta,
 \label{eq:An}\\
 B_n(x)&:=\frac1\pi\int_0^{2\pi}
             a_\theta(x)\sin(n\theta)\,d\theta,
 \label{eq:Bn}\\
 z_n(x)&:=A_n(x)+iB_n(x).
 \label{eq:zn}
\end{align}
Since $a_{\theta+\pi}=-a_\theta$, all even coefficients vanish.  The odd
coefficients agree with those of the original family $u_\theta$.

Let
\[
        F_n=(A_n,B_n):M\to\R^2.
\]

\begin{lemma}[Differentiating the coefficients]\label{lem:diff}
There is a full-measure subset of $M$ such that for every $x$ in this set,
for almost every $\theta$, $a_\theta$ is differentiable at $x$, and for
every $v\in T_xM$,
\begin{align*}
 dA_n(v)&=\frac1\pi\int_0^{2\pi}
               da_\theta(v)\cos(n\theta)\,d\theta,\\
 dB_n(v)&=\frac1\pi\int_0^{2\pi}
               da_\theta(v)\sin(n\theta)\,d\theta.
\end{align*}
\end{lemma}

\begin{proof}
This follows from joint Lipschitz regularity, Rademacher's theorem, Fubini,
and weak differentiation under the $\theta$-integral.
\end{proof}

We first state the usual unit Jacobian budget, and then its sharper slack
version.

\begin{proposition}[Orthogonal Jacobian budget]\label{prop:budget}
Let $m_1,\ldots,m_N$ be distinct positive odd integers and
$U=(u_{jk})\in O(N)$.  Put
\[
        G_j=\sum_{k=1}^Nu_{jk}F_{m_k}.
\]
Then almost everywhere,
\begin{equation}\label{eq:budget-unit}
        \sum_{j=1}^NJ_2G_j\le1.
\end{equation}
More precisely,
\begin{equation}\label{eq:budget-defect}
 \boxed{
 \sum_{j=1}^NJ_2G_j(x)
 \le
 1-\frac1{2\pi}\int_0^{2\pi}
       |\nabla\sigma_\theta(x)|^2\,d\theta.}
\end{equation}
\end{proposition}

\begin{proof}
Fix an orthonormal basis $(e_1,e_2)$ at a regular point and set
\[
        f(\theta)=da_\theta(e_1),\qquad
        g(\theta)=da_\theta(e_2).
\]
For the original Fourier maps set
\[
        p_k=dF_{m_k}(e_1),\qquad
        q_k=dF_{m_k}(e_2)\in\R^2.
\]
Bessel's inequality gives
\[
 \sum_k|p_k|^2\le\frac1\pi\int f^2,\qquad
 \sum_k|q_k|^2\le\frac1\pi\int g^2.
\]
Orthogonal mixing preserves these square sums.  Hence
\[
\begin{aligned}
 \sum_jJ_2G_j
 &\le\frac12\sum_j
     \bigl(|dG_j(e_1)|^2+|dG_j(e_2)|^2\bigr)\\
 &\le\frac1{2\pi}\int_0^{2\pi}
       \bigl(f(\theta)^2+g(\theta)^2\bigr)\,d\theta\\
 &=\frac1{2\pi}\int_0^{2\pi}|\nabla a_\theta|^2\,d\theta.
\end{aligned}
\]
Now use \eqref{eq:parallelogram-slack}.  Dropping the nonnegative defect
gives \eqref{eq:budget-unit}.
\end{proof}

\section{Exact boundary Fourier coefficients}

Let
\[
        q(t)=d_{\T}(t,0)-\frac{\pi}{2}.
\]
On $[-\pi,\pi]$, $q(t)=|t|-\pi/2$.

\begin{lemma}[Triangle-wave coefficients]\label{lem:triangle}
For $n\ge1$,
\[
 \frac1\pi\int_{-\pi}^{\pi}|t|\cos(nt)\,dt
 =
 \begin{cases}
 -\dfrac4{\pi n^2},&n\text{ odd},\\[1.5mm]
 0,&n\text{ even}.
 \end{cases}
\]
\end{lemma}

\begin{proof}
By evenness,
\[
 \frac2\pi\int_0^\pi t\cos(nt)\,dt
 =
 \frac{2}{\pi n^2}\bigl((-1)^n-1\bigr).
\]
\end{proof}

Consequently, for every positive odd $n$,
\begin{equation}\label{eq:boundary-zn}
        \boxed{
        z_n(\gamma(s))
        =-\rho_n e^{ins},
        \qquad
        \rho_n:=\frac4{\pi n^2}.}
\end{equation}

\section{The universal orientation-free bound}

We include the complete untwisting argument because it is what removes the
orientation hypothesis.

Fix $N$ and set
\[
        m_k=2k-1,\qquad w_k=\frac1{m_k},
        \qquad1\le k\le N.
\]
For $k\ge2$ define
\[
 c_k=(1+4w_k^2)^{-1/2},
 \qquad
 s_k=2w_k(1+4w_k^2)^{-1/2}.
\]
Let $R_k\in O(N)$ be the identity outside the $(1,k)$ coordinate plane and
have block
\[
        \begin{pmatrix}c_k&-s_k\\ s_k&c_k\end{pmatrix}
\]
on that plane.  Set
\[
        U_N=R_2R_3\cdots R_N=(u_{jk}).
\]

\begin{lemma}[Entries of the Givens product]\label{lem:Givens-entries}
The first row satisfies
\[
 u_{11}=\prod_{\ell=2}^Nc_\ell,\qquad
 u_{1k}=-s_k\prod_{\ell=2}^{k-1}c_\ell.
\]
For $j\ge2$,
\[
\begin{aligned}
 u_{j1}&=s_j\prod_{\ell=j+1}^Nc_\ell,\\
 u_{jk}&=0\quad(2\le k<j),\\
 u_{jj}&=c_j,\\
 u_{jk}&=-s_js_k\prod_{\ell=j+1}^{k-1}c_\ell
 \quad(k>j).
\end{aligned}
\]
\end{lemma}

\begin{proof}
Multiply successively on the right.  Before $R_k$ is inserted, column $k$
is still $e_k$.  The rotation replaces the first column $C_1$ and the
$k$th column by
\[
        c_kC_1+s_ke_k,\qquad -s_kC_1+c_ke_k.
\]
The formulas follow by induction.
\end{proof}

\begin{lemma}[First-harmonic dominance]\label{lem:dominance}
Every row of $U_N$ satisfies
\begin{equation}\label{eq:row-dominance}
        \boxed{
        |u_{j1}|>
        \sum_{k=2}^N\frac{|u_{jk}|}{m_k}.}
\end{equation}
\end{lemma}

\begin{proof}
We use
\[
 c_k=(1+4w_k^2)^{-1/2}
 \ge \exp(-2w_k^2).
\]
Hence, for any finite set $I$,
\[
        \prod_{k\in I}c_k
        \ge1-2\sum_{k\in I}w_k^2.
\]

For the first row, let $S=\sum_{k=2}^Nw_k^2$.  Then
\[
 \sum_{k=2}^Nw_k|u_{1k}|
 \le2S,
 \qquad
 |u_{11}|\ge1-2S.
\]
The strictly convex function $f(x)=(2x-1)^{-2}$ satisfies
\[
 \sum_{k=2}^{\infty}\frac1{(2k-1)^2}
 <
 \int_{3/2}^{\infty}\frac{dx}{(2x-1)^2}
 =\frac14.
\]
Thus $1-2S>2S$.

For $j\ge2$ let
\[
        S_j=\sum_{k=j+1}^Nw_k^2.
\]
The formulas of Lemma~\ref{lem:Givens-entries} give
\[
 |u_{j1}|\ge2w_jc_j(1-2S_j)
\]
and
\[
 \sum_{k=2}^Nw_k|u_{jk}|
 \le w_jc_j(1+4S_j).
\]
It suffices that $S_j<1/8$.  Since $j\ge2$,
\[
 \sum_{k=3}^{\infty}\frac1{(2k-1)^2}
 <
 \int_{5/2}^{\infty}\frac{dx}{(2x-1)^2}
 =\frac18.
\]
\end{proof}

\begin{lemma}[Dominant first harmonic]\label{lem:jordan}
Let
\[
        z(s)=\sum_{k=1}^Na_ke^{im_ks},
        \qquad1=m_1<m_2<\cdots<m_N,
\]
and assume
\[
        |a_1|>\sum_{k=2}^Nm_k|a_k|.
\]
Then $z:S^1\to\C$ is an embedding, has winding number $+1$ around $0$, and
its enclosed area is
\begin{equation}\label{eq:jordan-area}
        |\Omega_z|=\pi\sum_{k=1}^Nm_k|a_k|^2.
\end{equation}
\end{lemma}

\begin{proof}
For every positive integer $m$,
\[
        |e^{ims}-e^{imt}|
        \le m|e^{is}-e^{it}|.
\]
Therefore, if $s\ne t$,
\[
 |z(s)-z(t)|
 \ge
 \left(
 |a_1|-\sum_{k=2}^Nm_k|a_k|
 \right)|e^{is}-e^{it}|>0.
\]
Thus $z$ is a Jordan curve.  The higher-frequency part has modulus strictly
smaller than $|a_1|$, so straight-line homotopy to $a_1e^{is}$ avoids the
origin and the winding number is $+1$.  Green's formula and Fourier
orthogonality give \eqref{eq:jordan-area}.
\end{proof}

\begin{lemma}[Jordan coverage without orientation]\label{lem:mod2}
Let $M$ be any compact connected surface with one boundary component and
$G:M\to\R^2$ continuous.  If $G|_{\partial M}$ is a Jordan curve $C$, then
the bounded component $\Omega_C$ of $\R^2\setminus C$ lies in $G(M)$.
Consequently, if $G$ is Lipschitz,
\[
        \int_MJ_2G\,dA\ge|\Omega_C|.
\]
\end{lemma}

\begin{proof}
If $y\in\Omega_C\setminus G(M)$, radial projection around $y$ gives a map
$M\to S^1$ whose restriction to $\partial M$ has odd degree.  This is
impossible over $\Z_2$: the boundary class maps to zero in
$H_1(M;\Z_2)$ because it is the boundary of the relative fundamental class
$[M,\partial M]_2$.  The area conclusion is the ordinary two-dimensional
area formula.
\end{proof}

We now prove the finite certificate.  Define
\[
        G_j=\sum_{k=1}^Nu_{jk}F_{m_k}.
\]
By \eqref{eq:boundary-zn}, its boundary is, up to an overall sign,
\[
        \sum_{k=1}^Nu_{jk}\rho_{m_k}e^{im_ks}.
\]
Since
\[
        m_k\rho_{m_k}
        =\rho_1\,\frac1{m_k},
\]
Lemma~\ref{lem:dominance} implies the hypothesis of
Lemma~\ref{lem:jordan} for every row.

\begin{proposition}[Finite universal certificate]\label{prop:finite-univ}
For every $N\ge1$,
\[
        \boxed{
        \Area(M)\ge
        \frac{16}{\pi}
        \sum_{k=1}^N\frac1{(2k-1)^3}.}
\]
No orientability assumption is required.
\end{proposition}

\begin{proof}
Lemma~\ref{lem:jordan} gives
\[
 |\Omega_j|
 =
 \pi\sum_{k=1}^Nm_k\rho_{m_k}^2u_{jk}^2.
\]
By Lemma~\ref{lem:mod2} and Proposition~\ref{prop:budget},
\[
\begin{aligned}
 \Area(M)
 &\ge\sum_{j=1}^N\int_MJ_2G_j\,dA\\
 &\ge\sum_{j=1}^N|\Omega_j|\\
 &=\pi\sum_{k=1}^Nm_k\rho_{m_k}^2
   \sum_{j=1}^Nu_{jk}^2\\
 &=\frac{16}{\pi}
   \sum_{k=1}^N\frac1{m_k^3}.
\end{aligned}
\]
\end{proof}

Letting $N\to\infty$ and using
\[
 \sum_{k=1}^{\infty}\frac1{(2k-1)^3}
 =\left(1-\frac18\right)\zeta(3)
 =\frac78\zeta(3)
\]
proves Theorem~\ref{thm:universal}.

\section{An antipodal-slack defect refinement}

Keeping the sharper budget \eqref{eq:budget-defect} instead of discarding
the slack gives a useful nonlinear refinement.

\begin{theorem}[Antipodal defect]\label{thm:slack}
Every compact Riemannian isometric filling, orientable or not, satisfies
\begin{equation}\label{eq:slack-bound}
 \boxed{
 \Area(M)\ge\frac{14\zeta(3)}{\pi}
 +\frac1{2\pi}\int_M\int_0^{2\pi}
          |\nabla\sigma_\theta(x)|^2\,d\theta\,dA(x),}
\end{equation}
where
\[
        \sigma_\theta
        =\frac{u_\theta+u_{\theta+\pi}-\pi}{2}\ge0.
\]
\end{theorem}

\begin{proof}
Use \eqref{eq:budget-defect} in the proof of
Proposition~\ref{prop:finite-univ}.  It yields
\[
 \Area(M)
 -
 \frac1{2\pi}
 \int_M\int_0^{2\pi}|\nabla\sigma_\theta|^2
 \ge
 \frac{16}{\pi}\sum_{k=1}^N\frac1{(2k-1)^3}.
\]
Let $N\to\infty$.
\end{proof}

\begin{remark}
The defect vanishes exactly when the antipodal equalities
\[
        u_\theta+u_{\theta+\pi}=\pi
\]
hold infinitesimally.  These equalities characterize the anti-periodic
distance-profile geometry underlying the injective hull.  Thus
Theorem~\ref{thm:slack} quantifies how much area is lost when one projects a
general boundary-distance profile to its anti-periodic part.
\end{remark}

\section{The oriented linear calibration}

For the nonlinear analysis it is convenient to return to oriented tangent
planes and complex notation.

Let $\mathcal O=\{1,3,5,\ldots\}$ and
\[
        \mathcal H=\ell^2(\mathcal O;\C).
\]
The Fourier coefficient map is
\[
        Z:M\to\mathcal H,\qquad
        Z(x)=(z_n(x))_{n\in\mathcal O}.
\]
Bessel's inequality shows that $Z$ is Lipschitz.

On $\mathcal H$ consider the standard symplectic form
\[
        \omega(p,q)
        =\sum_{n\in\mathcal O}\operatorname{Im}
          (\overline{p_n}q_n).
\]
It has primitive
\[
        \alpha_z(v)=\frac12\omega(z,v).
\]

\begin{lemma}[Hilbert-valued Stokes]\label{lem:Hilbert-Stokes}
The map $Z:M\to\mathcal H$ is Lipschitz.  More generally, if
$G:M\to\mathcal H$ is a Lipschitz map, then
\[
        \int_M G^*\omega=\int_{\partial M}G^*\alpha.
\]
\end{lemma}

\begin{proof}
If $x,y\in M$, then
$|a_\theta(x)-a_\theta(y)|\le d_M(x,y)$.  Bessel's inequality therefore
implies
\[
 \|Z(x)-Z(y)\|_{\ell^2}^2
 \le\frac1\pi\int_0^{2\pi}
       |a_\theta(x)-a_\theta(y)|^2\,d\theta
 \le2d_M(x,y)^2.
\]
Thus $Z$ is Lipschitz.

For the Stokes identity let $P_N$ be projection onto the first $N$ odd
coordinates.  Ordinary weak Stokes applied componentwise gives
\[
 \int_M(P_NG)^*\omega
 =\int_{\partial M}(P_NG)^*\alpha.
\]
At almost every point $DG$ is a bounded linear map into $\mathcal H$, and
$P_NDG\to DG$ strongly.  The symplectic integrands converge pointwise and
are dominated by
$\|DG(e_1)\|\,\|DG(e_2)\|$, which is integrable for a Lipschitz map on the
compact surface.  The boundary primitives converge in the same way.  Let
$N\to\infty$.
\end{proof}

Fix an oriented orthonormal frame $(e_1,e_2)$ and set
\[
        W(\theta)
        =da_\theta(e_1)+i\,da_\theta(e_2).
\]
Write
\[
        W(\theta)=\sum_{k\in\Z}c_ke^{ik\theta}
\]
with the standard normalization
\[
        c_k=\frac1{2\pi}\int_0^{2\pi}
             W(\theta)e^{-ik\theta}\,d\theta.
\]
Because $a_{\theta+\pi}=-a_\theta$, only odd $k$ occur.  Since
$|W(\theta)|\le1$,
\[
        \sum_k|c_k|^2\le1.
\]

A direct calculation gives
\begin{equation}\label{eq:base-density}
 \boxed{
 (Z^*\omega)(e_1,e_2)
 =
 \sum_{n\in\mathcal O}
 \bigl(|c_n|^2-|c_{-n}|^2\bigr).}
\end{equation}
Thus the comass is at most one.  Equality $+1$ occurs precisely when
$|W|=1$ almost everywhere and all negative Fourier coefficients vanish;
equality $-1$ is the anti-analytic analogue.

On the boundary,
\[
        z_n(\gamma(s))=-\rho_ne^{ins}.
\]
Therefore
\[
\begin{aligned}
 \int_{\partial M}Z^*\alpha
 &=
 \pi\sum_{n\in\mathcal O}n\rho_n^2\\
 &=
 \frac{16}{\pi}\sum_{n\in\mathcal O}\frac1{n^3}\\
 &=
 \frac{14\zeta(3)}{\pi}.
\end{aligned}
\]
This recovers the oriented linear bound.

\section{A barrier for separable nonlinear reparametrizations}

One natural attempt is to replace the scalar distance $u_\theta$ by
$\psi(u_\theta)$ for a $1$-Lipschitz scalar function $\psi$ and then repeat
the Fourier calibration.  This cannot improve the linear constant.

\begin{proposition}[Separable nonlinear barrier]\label{prop:separable}
Let $\psi:[0,\pi]\to\R$ be absolutely continuous with
$|\psi'|\le1$ almost everywhere.  Apply the oriented Fourier calibration
to the family $\psi(u_\theta)$.  Its boundary action is at most
\[
        \frac{14\zeta(3)}{\pi}.
\]
The same conclusion holds for a Hilbert-valued collection of scalar
features whose pointwise derivative vector has norm at most one.
\end{proposition}

\begin{proof}
Put $p=\psi'$.  The $n$th cosine coefficient of the boundary function
$\psi(d_{\T}(\theta,s))$ has radius
\[
 r_n^\psi
 =
 \frac{2}{\pi n}
 \left|
 \int_0^\pi p(t)\sin(nt)\,dt
 \right|.
\]
Hence the total oriented boundary action is
\begin{equation}\label{eq:Qp}
 \mathcal Q(p)
 =
 \frac4\pi\sum_{n\ge1}\frac1n
 \left(
 \int_0^\pi p(t)\sin(nt)\,dt
 \right)^2.
\end{equation}
The kernel of this quadratic form is
\begin{equation}\label{eq:positive-kernel}
 K(t,s)
 =
 \sum_{n\ge1}\frac{\sin(nt)\sin(ns)}{n}
 =
 \frac12\log
 \frac{\sin((t+s)/2)}
      {\sin(|t-s|/2)}.
\end{equation}
For $0<t,s<\pi$ this kernel is positive, because
\[
 \sin^2\frac{t+s}{2}
 -
 \sin^2\frac{|t-s|}{2}
 =
 \sin t\,\sin s>0.
\]
Therefore, under $|p|\le1$,
\[
 \mathcal Q(p)
 \le
 \frac4\pi\int_0^\pi\int_0^\pi K(t,s)\,dt\,ds
 =
 \mathcal Q(1).
\]
The choice $p\equiv1$ is the original distance function and gives
$14\zeta(3)/\pi$.

For a Hilbert-valued derivative $p(t)$ with $|p(t)|\le1$, replace
$p(t)p(s)$ by $\langle p(t),p(s)\rangle$ and use
$\langle p(t),p(s)\rangle\le|p(t)||p(s)|\le1$.
\end{proof}

Thus any improvement has to couple distinct boundary parameters or Fourier
modes genuinely nonlinearly.

\section{A Euclidean Hilbert-space barrier}

There is a second useful obstruction.  Consider the boundary Fourier loop
in $\mathcal H$,
\[
        \Gamma(s)=(-\rho_ne^{ins})_{n\in\mathcal O}.
\]
It spans the holomorphic disk
\begin{equation}\label{eq:holodisk}
        \Phi(\zeta)=(-\rho_n\zeta^n)_{n\in\mathcal O},
        \qquad|\zeta|\le1.
\end{equation}

\begin{proposition}[Holomorphic-disk barrier]\label{prop:holodisk}
The ordinary Euclidean area of the disk \eqref{eq:holodisk} in the Hilbert
space $\mathcal H$ is exactly
\[
        \frac{14\zeta(3)}{\pi}.
\]
Consequently, no closed two-form on $\mathcal H$ with ordinary Euclidean
comass at most one can have larger action on $\Gamma$.
\end{proposition}

\begin{proof}
The holomorphic map $\Phi$ is conformal, and
\[
 \|\Phi'(\zeta)\|_{\ell^2}^2
 =
 \sum_{n\in\mathcal O}n^2\rho_n^2|\zeta|^{2n-2}.
\]
Therefore
\[
\begin{aligned}
 \Area_{\rm Euc}(\Phi)
 &=
 2\pi\int_0^1
 \sum_{n\in\mathcal O}n^2\rho_n^2r^{2n-2}\,r\,dr\\
 &=
 \pi\sum_{n\in\mathcal O}n\rho_n^2
 =
 \frac{14\zeta(3)}{\pi}.
\end{aligned}
\]
Stokes and the comass inequality give the second statement.
\end{proof}

\begin{remark}
The nonlinear improvement below does not contradict
Proposition~\ref{prop:holodisk}: its comass is controlled only on the much
smaller class of tangent planes arising from pointwise-bounded
distance-profile derivatives.  The special $L^\infty$ geometry, not the
Euclidean Hilbert geometry, is what supplies the gain.
\end{remark}

\section{A phase-resonant nonlinear deformation}

We now assume $M$ is oriented.

For $\lambda\ge0$ define the nonlinear map
\begin{equation}\label{eq:G-lambda}
 \mathcal G_\lambda:\mathcal H\to\mathcal H,\qquad
 (\mathcal G_\lambda(z))_n
 =
 z_n+\lambda\,\overline{z_1}^{\,2}z_{n+2},
 \qquad n\in\mathcal O.
\end{equation}
Let
\[
        \Omega_\lambda
        :=
        (\mathcal G_\lambda\circ Z)^*\omega.
\]
The form is exact.  One may justify Stokes by projecting to finitely many
coordinates and passing to the limit; all relevant derivative sequences
are square summable.

The reason for \eqref{eq:G-lambda} is the boundary resonance:
\[
 \overline{z_1(\gamma(s))}^{\,2}z_{n+2}(\gamma(s))
 =
 -\rho_1^2\rho_{n+2}e^{ins}.
\]
Thus
\begin{equation}\label{eq:G-boundary}
 (\mathcal G_\lambda\circ Z)_n(\gamma(s))
 =
 -\bigl(\rho_n+\lambda\rho_1^2\rho_{n+2}\bigr)e^{ins}.
\end{equation}

\subsection{Exact boundary action}

Let
\[
        \mathcal B(\lambda)
        :=
        \int_{\partial M}
        (\mathcal G_\lambda\circ Z)^*\alpha.
\]
By \eqref{eq:G-boundary},
\begin{equation}\label{eq:B-series}
 \mathcal B(\lambda)
 =
 \pi\sum_{n\in\mathcal O}
 n\bigl(\rho_n+\lambda\rho_1^2\rho_{n+2}\bigr)^2.
\end{equation}

\begin{proposition}[Boundary action]\label{prop:B}
One has
\begin{equation}\label{eq:Bclosed}
 \boxed{
 \mathcal B(\lambda)
 =
 C_0+b\lambda+d\lambda^2,}
\end{equation}
where
\begin{align}
 C_0&=\frac{14\zeta(3)}{\pi},
 \label{eq:C0}\\
 b&=
 \frac{512}{\pi^3}
 \left(\frac34-\frac{\pi^2}{16}\right)
 =
 2.1986728644\ldots,
 \label{eq:b}\\
 d&=
 \frac{4096}{\pi^5}
 \left(
 \frac78\zeta(3)+1-\frac{\pi^4}{48}
 \right)
 =
 0.3004038975\ldots.
 \label{eq:d}
\end{align}
In particular, $b>0$.
\end{proposition}

\begin{proof}
The constant term is the linear Fourier action.  For the linear term,
\[
 \sum_{n\in\mathcal O}\frac1{n(n+2)^2}
 =
 \frac34-\frac{\pi^2}{16},
\]
which follows from the partial fraction decomposition
\[
 \frac1{n(n+2)^2}
 =
 \frac1{4n}-\frac1{4(n+2)}-\frac1{2(n+2)^2}.
\]
For the quadratic term,
\[
\begin{aligned}
 \sum_{n\in\mathcal O}\frac{n}{(n+2)^4}
 &=
 \sum_{\substack{m\ge3\\m\ {\rm odd}}}
 \left(\frac1{m^3}-\frac2{m^4}\right)\\
 &=
 \frac78\zeta(3)+1-\frac{\pi^4}{48}.
\end{aligned}
\]
Substitution into \eqref{eq:B-series} gives the formulas.
\end{proof}

\section{The comass cancellation}

Fix a regular point and an oriented orthonormal frame $(e_1,e_2)$.  Write
\[
 p_n=dz_n(e_1),\qquad q_n=dz_n(e_2).
\]
Let
\[
 W(\theta)
 =da_\theta(e_1)+i\,da_\theta(e_2)
 =\sum_{k\ {\rm odd}}c_ke^{ik\theta}.
\]
Then
\begin{equation}\label{eq:pq-c}
 p_n=c_{-n}+\overline{c_n},
 \qquad
 q_n=i(\overline{c_n}-c_{-n}).
\end{equation}
Consequently
\[
        \Omega_0(e_1,e_2)
        =
        \sum_{n\in\mathcal O}
        (|c_n|^2-|c_{-n}|^2).
\]

Expand
\[
        \Omega_\lambda
        =
        \Omega_0+\lambda L+\lambda^2Q.
\]

Let
\[
        a=z_1,\qquad y_n=z_{n+2}.
\]
A direct calculation from \eqref{eq:G-lambda} gives the following exact
formula.

\begin{lemma}[Exact first variation]\label{lem:firstvar}
At every regular oriented tangent plane,
\[
\begin{aligned}
 L
 &=
 2\operatorname{Re}\left[
 \overline a^{\,2}
 \sum_{n\in\mathcal O}
 \left(
 c_n\overline{c_{n+2}}
 -
 \overline{c_{-n}}c_{-(n+2)}
 \right)
 \right]\\
 &\quad+
 4\operatorname{Re}\left[
 \overline a
 \sum_{n\in\mathcal O}
 y_n
 \left(
 -\overline{c_{-n}}c_1
 +\overline{c_{-1}}c_n
 \right)
 \right].
\end{aligned}
\]
\end{lemma}

\begin{proof}
Differentiate
\[
        \overline{z_1}^{\,2}z_{n+2}.
\]
For a tangent vector with coefficient sequence $p$,
\[
        r_n(p)
        =
        2\overline a\,y_n\,\overline{p_1}
        +\overline a^{\,2}p_{n+2}.
\]
The first variation of the $n$th symplectic Jacobian is
\[
 \operatorname{Im}
 \bigl(\overline{r_n(p)}q_n+\overline{p_n}r_n(q)\bigr).
\]
For later reference, \eqref{eq:pq-c} gives the elementary identities
\[
 \frac1{2i}(d\overline z_j\wedge dz_k)(e_1,e_2)
 =c_j\overline{c_k}-\overline{c_{-j}}c_{-k},
\]
\[
 \frac1{2i}(d\overline z_j\wedge d\overline z_k)(e_1,e_2)
 =-\overline{c_{-j}}c_k+\overline{c_{-k}}c_j.
\]
Expanding
$d(\overline z_1^{\,2}z_{n+2})$ and pairing the two conjugate terms gives
exactly the displayed formula for $L$.
\end{proof}

If $\Omega_0=1$, then $c_{-n}=0$, $|W|=1$, and
\[
 \sum_{n\in\mathcal O}c_n\overline{c_{n+2}}=0
\]
because it is the nonzero Fourier coefficient of $|W|^2\equiv1$.
Therefore $L=0$.  The same holds at $\Omega_0=-1$.

\begin{corollary}[First-order stationarity]\label{cor:stationary}
The first variation $L$ vanishes on every plane at which
$|\Omega_0|=1$.
\end{corollary}

This is the nonlinear mechanism: the boundary action changes to first
order, while the comass does not.

\section{Quantitative comass estimate}

We now quantify Corollary~\ref{cor:stationary} with explicit constants.

\subsection{Profile coefficient bounds}

Because $\theta\mapsto a_\theta(x)$ is $1$-Lipschitz and periodic,
\[
        \sum_{n\ge1}n^2|z_n|^2
        \le2.
\]
Moreover,
\begin{equation}\label{eq:z1-bound}
        |z_1|\le\frac4\pi.
\end{equation}
Indeed, integration by parts gives
\[
 z_1=-\frac1{i\pi}\int_0^{2\pi}
          \partial_\theta a_\theta\,e^{i\theta}\,d\theta.
\]
After rotating the complex phase,
\[
 \left|
 \int \partial_\theta a_\theta\,e^{i\theta}
 \right|
 \le\int_0^{2\pi}|\cos\theta|\,d\theta=4.
\]

Set
\[
        r=|z_1|,
        \qquad
        Y^2=\sum_{n\in\mathcal O}|z_{n+2}|^2.
\]
Then
\begin{equation}\label{eq:Y-bound}
        \boxed{
        Y^2\le\frac{2-r^2}{9}.}
\end{equation}

\subsection{Fourier deficit}

Assume first that $\Omega_0\ge0$ and put
\[
 P=\sum_{n\in\mathcal O}|c_n|^2,\qquad
 N=\sum_{n\in\mathcal O}|c_{-n}|^2,
\]
\[
        \delta=1-\Omega_0=1-P+N.
\]
Since $P+N\le1$,
\begin{equation}\label{eq:Ndelta}
        N\le\frac{\delta}{2}.
\end{equation}
Also
\[
        1-P-N=\delta-2N.
\]

Because $1-|W|^2\ge0$, every nonzero Fourier coefficient of $|W|^2$ has
absolute value at most its mean deficit:
\[
 \left|
 \sum_{k\in\Z}c_k\overline{c_{k+2}}
 \right|
 \le1-P-N=\delta-2N.
\]
Set
\[
 A:=\sum_{n\in\mathcal O}c_n\overline{c_{n+2}},\qquad
 B:=\sum_{n\in\mathcal O}\overline{c_{-n}}c_{-(n+2)},\qquad
 X:=c_{-1}\overline{c_1}.
\]
Since only odd coefficients occur, the full shift-two correlation is
\[
 \sum_{k\in\Z}c_k\overline{c_{k+2}}=A+B+X.
\]
Moreover $|B|\le N$ by Cauchy--Schwarz and
$|X|\le\sqrt N$.  Hence
\[
 |A-B|
 =|(A+B+X)-X-2B|
 \le(\delta-2N)+\sqrt N+2N
 =\delta+\sqrt N.
\]
Using $N\le\delta/2$ gives
\begin{equation}\label{eq:Bcorr}
 \boxed{
 \left|
 \sum_{n\in\mathcal O}
 \left(
 c_n\overline{c_{n+2}}
 -
 \overline{c_{-n}}c_{-(n+2)}
 \right)
 \right|
 \le\delta+\sqrt{\delta/2}.}
\end{equation}
Furthermore, Cauchy--Schwarz gives
\begin{equation}\label{eq:ycorr}
 \left|
 \sum_{n\in\mathcal O}
 y_n
 \left(
 -\overline{c_{-n}}c_1
 +\overline{c_{-1}}c_n
 \right)
 \right|
 \le2Y\sqrt N
 \le\sqrt2\,Y\sqrt\delta.
\end{equation}

Combining Lemma~\ref{lem:firstvar},
\eqref{eq:Bcorr}, and \eqref{eq:ycorr} yields
\begin{equation}\label{eq:L-r}
 |L|
 \le
 2r^2\delta
 +
 \sqrt2\bigl(r^2+4rY\bigr)\sqrt\delta.
\end{equation}
Using \eqref{eq:Y-bound},
\[
 \sqrt2(r^2+4rY)
 \le
 \sqrt2\left[
 r^2+\frac43r\sqrt{2-r^2}
 \right].
\]
Elementary calculus shows that the maximum occurs at $r^2=8/5$ and equals
$8\sqrt2/3$; the critical point lies in the allowed interval because
$\pi^2<10$.  Since $r\le4/\pi$, we obtain
\begin{equation}\label{eq:Lbound}
 \boxed{
 |L|
 \le C_*\sqrt\delta+D_*\delta,}
\end{equation}
where
\begin{equation}\label{eq:CD}
        C_*=\frac{8\sqrt2}{3},
        \qquad
        D_*=\frac{32}{\pi^2}.
\end{equation}
The case $\Omega_0\le0$ is identical after reversing the oriented frame.
Hence \eqref{eq:Lbound} holds with
\[
        \delta=1-|\Omega_0|.
\]

\subsection{Quadratic term}

For a derivative sequence $p=(p_n)_{n\in\mathcal O}$, write
\[
        R(p)_n
        =
        2\overline a\,y_n\,\overline{p_1}
        +\overline a^{\,2}p_{n+2}.
\]
If
\[
        x=|p_1|,\qquad
        y=\left(\sum_{n\ge3,\ n\ {\rm odd}}|p_n|^2\right)^{1/2},
\]
then
\[
 \|R(p)\|_{\ell^2}
 \le2rYx+r^2y.
\]
Therefore
\[
 \|R(p)\|_{\ell^2}^2
 \le
 \bigl(r^4+4r^2Y^2\bigr)\|p\|_{\ell^2}^2.
\]
By \eqref{eq:Y-bound},
\[
 r^4+4r^2Y^2
 \le
 \frac59r^4+\frac89r^2.
\]
The right side is increasing in $r^2$, so using $r\le4/\pi$ gives
\begin{equation}\label{eq:Qstar}
        \boxed{
        |Q|\le Q_*,
        \qquad
        Q_*=
        \frac{128}{9\pi^2}
        +\frac{1280}{9\pi^4}.}
\end{equation}

Indeed,
\[
 |Q|
 \le\|R(p)\|_2\|R(q)\|_2
 \le\frac{Q_*}{2}
       \bigl(\|p\|_2^2+\|q\|_2^2\bigr)
 \le Q_*,
\]
because Bessel gives
$\|p\|_2^2+\|q\|_2^2\le2$.

\subsection{Global comass}

\begin{proposition}[Nonlinear comass bound]\label{prop:comass}
For
\[
        0\le\lambda<\frac{\pi^2}{32},
\]
the nonlinear form satisfies
\begin{equation}\label{eq:comass}
 \boxed{
 \|\Omega_\lambda\|_{\rm comass}
 \le
 \mathcal C(\lambda):=
 1+\lambda^2
 \left[
 Q_*+
 \frac{C_*^2}{4(1-D_*\lambda)}
 \right].}
\end{equation}
\end{proposition}

\begin{proof}
At an arbitrary oriented orthonormal plane let
\[
        \delta=1-|\Omega_0|\in[0,1].
\]
By \eqref{eq:Lbound} and \eqref{eq:Qstar},
\[
 |\Omega_\lambda|
 \le
 1-\delta
 +\lambda(C_*\sqrt\delta+D_*\delta)
 +Q_*\lambda^2.
\]
Thus
\[
 |\Omega_\lambda|
 \le
 1
 -(1-D_*\lambda)\delta
 +C_*\lambda\sqrt\delta
 +Q_*\lambda^2.
\]
For $A>0$,
\[
        -A\delta+B\sqrt\delta
        \le\frac{B^2}{4A}.
\]
Take
$A=1-D_*\lambda$ and $B=C_*\lambda$.
\end{proof}

\section{The explicit nonlinear lower bound}

Stokes gives
\[
        \int_M\Omega_\lambda=\mathcal B(\lambda).
\]
Combining Propositions~\ref{prop:B} and \ref{prop:comass} proves the
following.

\begin{theorem}[One-parameter nonlinear certificate]
\label{thm:nonlinear-formula}
For every compact oriented Riemannian isometric filling and every
\[
        0\le\lambda<\frac{\pi^2}{32},
\]
\begin{equation}\label{eq:nonlinear-master}
 \boxed{
 \Area(M)\ge
 \frac{
 C_0+b\lambda+d\lambda^2
 }{
 1+\lambda^2
 \left[
 Q_*+
 \dfrac{C_*^2}{4(1-D_*\lambda)}
 \right]
 }.}
\end{equation}
The constants are
\[
\begin{gathered}
 C_0=\frac{14\zeta(3)}{\pi},\qquad
 b=\frac{512}{\pi^3}
 \left(\frac34-\frac{\pi^2}{16}\right),\\
 d=\frac{4096}{\pi^5}
 \left(\frac78\zeta(3)+1-\frac{\pi^4}{48}\right),\\
 C_*=\frac{8\sqrt2}{3},\qquad
 D_*=\frac{32}{\pi^2},\qquad
 Q_*=\frac{128}{9\pi^2}+\frac{1280}{9\pi^4}.
\end{gathered}
\]
\end{theorem}

For $\lambda=0.03$,
\[
 \mathcal B(0.03)
 =5.4230028938\ldots,
\]
\[
 \mathcal C(0.03)
 =1.0061557533\ldots,
\]
and therefore
\[
        \Area(M)
        \ge5.3898244639\ldots.
\]
This proves Theorem~\ref{thm:nonlinear-main}.  Numerically optimizing the
explicit right side of \eqref{eq:nonlinear-master} gives
\[
        \lambda\approx0.029227,
        \qquad
        \Area(M)\ge5.3898489672\ldots.
\]

\begin{remark}
The gain is not merely a numerical accident.  The numerator has positive
first derivative $b$ at $\lambda=0$, while the comass upper bound has zero
first derivative.  Hence every sufficiently small positive $\lambda$
strictly improves the linear Fourier constant.
\end{remark}

\section{The full resonant hierarchy}

The cubic resonance is only the first member of a natural family.  For
$j\ge1$ define
\begin{equation}\label{eq:Tj}
        (T_jz)_n
        =
        \overline{z_1}^{\,2j}z_{n+2j}.
\end{equation}
On the boundary,
\[
        (T_jZ(\gamma(s)))_n
        =
        -\rho_1^{2j}\rho_{n+2j}e^{ins}.
\]
Thus every $T_j$ increases the $n$th boundary radius coherently.

Consider
\[
        \mathcal G_{\boldsymbol\lambda}(z)
        =
        z+\sum_{j=1}^J\lambda_jT_jz.
\]
Its boundary action has a positive linear part
\[
        C_0+\sum_{j=1}^Jb_j\lambda_j+O(|\boldsymbol\lambda|^2),
\]
where
\begin{equation}\label{eq:bj}
 b_j=
 \frac{32}{\pi}
 \left(\frac{16}{\pi^2}\right)^j
 \sum_{n\in\mathcal O}
 \frac1{n(n+2j)^2}>0.
\end{equation}

At a positive saturation plane,
$W$ is an analytic inner function, so for every $j\ge1$,
\begin{equation}\label{eq:all-autocorr}
        \sum_{n\in\mathcal O}
        c_n\overline{c_{n+2j}}=0.
\end{equation}
At a negative saturation plane the analogous anti-analytic identity holds.
The same calculation as in Lemma~\ref{lem:firstvar} gives:

\begin{proposition}[Infinite family of stationary resonances]
\label{prop:hierarchy}
For every finite linear combination
$\sum_j\lambda_jT_j$, the first variation of the Fourier comass vanishes
on the entire saturation set $|\Omega_0|=1$, while the boundary action has
strictly positive first variation in every nonzero direction with
$\lambda_j\ge0$.
\end{proposition}

Thus the nonlinear problem has an infinite-dimensional cone of directions
with
\[
        \text{boundary gain}=O(|\boldsymbol\lambda|),
        \qquad
        \text{comass cost}=O(|\boldsymbol\lambda|^2).
\]
The next quantitative problem is to compute the true second-variation
quadratic form on this cone.  The crude independent estimates used above
suggest that the first shift $j=1$ is the most efficient single resonance,
but they do not resolve the optimal coupled problem.


\section{A rigorous ceiling for one-high-leg resonances}
\label{sec:one-high-leg-ceiling}

The preceding hierarchy raises a natural quantitative question: can one keep
adding the resonances $\overline z_1^{\,2j}z_{n+2j}$ and eventually approach the
hemisphere constant $2\pi$?  In this section we show that the answer is no for a
precise and fairly large class of such deformations.  In fact, the entire
normalized one-high-leg hierarchy has a rigorous ceiling below $5.466$.

Set
\[
        R:=\rho_1=\frac4\pi.
\]
For $J\ge1$ and $\boldsymbol\mu=(\mu_1,\ldots,\mu_J)\in\R^J$, consider the
normalized deformation
\begin{equation}\label{eq:normalized-one-leg}
 (\mathcal G_{\boldsymbol\mu}(z))_n
 =z_n+\sum_{j=1}^J
 \mu_j\left(\frac{\overline z_1}{R}\right)^{2j}z_{n+2j},
 \qquad n\in\mathcal O.
\end{equation}
On the boundary, the normalizing factor has unit modulus and exactly removes
$2j$ units of phase.  Hence
\begin{equation}\label{eq:normalized-boundary}
 (\mathcal G_{\boldsymbol\mu}(Z(\gamma(s))))_n
 =-\left(\rho_n+\sum_{j=1}^J\mu_j\rho_{n+2j}\right)e^{ins}.
\end{equation}

The point of the normalization is that the unavoidable high-frequency comass
cost becomes exactly the Euclidean square norm of the parameter vector.

\begin{lemma}[High-mode test]\label{lem:high-mode-test}
Let $\mathfrak C_J(\boldsymbol\mu)$ be any global ambient comass constant for the
pullback of the standard symplectic form under
$\mathcal G_{\boldsymbol\mu}$, where tangent planes are tested against all
anti-periodic fields
\[
        W(\theta)=v(\theta)+iw(\theta),
        \qquad |W(\theta)|\le1.
\]
Then
\begin{equation}\label{eq:high-mode-cost}
        \boxed{
        \mathfrak C_J(\boldsymbol\mu)
        \ge1+\sum_{j=1}^J\mu_j^2.}
\end{equation}
\end{lemma}

\begin{proof}
Evaluate at the boundary coefficient profile $Z(\gamma(0))$ and choose an odd
integer $m>2J$.  Take the unit-modulus positive Hardy mode
\[
        W_m(\theta)=e^{im\theta}.
\]
For the corresponding oriented plane, the Fourier derivative sequences satisfy
$p_m=1$, $q_m=i$, and all other entries vanish.  Since $m>2J$, the derivative
of the low-frequency factor $(\overline z_1/R)^{2j}$ vanishes on this plane,
while the shifted term $z_{n+2j}$ contributes precisely in coordinate
$n=m-2j$.  Thus
\[
 D\mathcal G_{\boldsymbol\mu}(p)
 =e_m+\sum_{j=1}^J\mu_je_{m-2j},
 \qquad
 D\mathcal G_{\boldsymbol\mu}(q)
 =iD\mathcal G_{\boldsymbol\mu}(p).
\]
The output coordinates are mutually orthogonal, so the standard symplectic form
evaluates to
\[
        1+\sum_{j=1}^J\mu_j^2.
\]
This plane is admissible because $|W_m|=1$.
\end{proof}

\subsection{The resonance matrix}

Put $x_0=1$ and $x_j=\mu_j$ for $1\le j\le J$.  The boundary action of
\eqref{eq:normalized-one-leg} is
\begin{equation}\label{eq:one-leg-boundary-action}
 \mathcal B_J(\boldsymbol\mu)
 =\pi\sum_{n\in\mathcal O}n
 \left(\sum_{j=0}^Jx_j\rho_{n+2j}\right)^2
 =x^TM^{(J)}x,
\end{equation}
where $M^{(J)}=(M_{jk})_{0\le j,k\le J}$ and
\begin{equation}\label{eq:resonance-matrix}
 \boxed{
 M_{jk}
 =\frac{16}{\pi}
 \sum_{n\in\mathcal O}
 \frac{n}{(n+2j)^2(n+2k)^2}.}
\end{equation}
The matrix is positive because it is a Gram matrix.

\begin{theorem}[Finite one-high-leg ceiling]\label{thm:finite-one-leg-ceiling}
For every $J$ and every parameter vector $\boldsymbol\mu$, any filling-area
certificate obtained from \eqref{eq:normalized-one-leg} by a global ambient
comass estimate satisfies
\begin{equation}\label{eq:finite-rayleigh-ceiling}
 \boxed{
 \frac{\mathcal B_J(\boldsymbol\mu)}
      {\mathfrak C_J(\boldsymbol\mu)}
 \le\lambda_{\max}(M^{(J)}).}
\end{equation}
\end{theorem}

\begin{proof}
By Lemma~\ref{lem:high-mode-test},
\[
 \frac{\mathcal B_J(\boldsymbol\mu)}
      {\mathfrak C_J(\boldsymbol\mu)}
 \le
 \frac{x^TM^{(J)}x}{x^Tx},
\]
and the right side is bounded by the largest eigenvalue.
\end{proof}

Let $M=(M_{jk})_{j,k\ge0}$ denote the infinite matrix on
$\ell^2(\mathbb N_0)$.  Positivity and the following trace computation show that
it is trace class.

\begin{lemma}[Exact trace]\label{lem:resonance-trace}
The positive operator $M$ is trace class and
\begin{equation}\label{eq:trace-M}
 \boxed{
 \operatorname{tr}M
 =\frac\pi2+\frac{7\zeta(3)}\pi+\frac{\pi^3}{24}.}
\end{equation}
\end{lemma}

\begin{proof}
By Tonelli,
\[
 \operatorname{tr}M
 =\frac{16}{\pi}
 \sum_{j\ge0}\sum_{n\in\mathcal O}
 \frac{n}{(n+2j)^4}.
\]
Write $m=n+2j$.  For fixed positive odd $m$, the possible values of $n$ are
$1,3,\ldots,m$, whose sum is $((m+1)/2)^2$.  Therefore
\[
\begin{aligned}
 \operatorname{tr}M
 &=\frac4\pi\sum_{m\in\mathcal O}
 \frac{(m+1)^2}{m^4}\\
 &=\frac4\pi\left(
 \sum_{m\in\mathcal O}\frac1{m^2}
 +2\sum_{m\in\mathcal O}\frac1{m^3}
 +\sum_{m\in\mathcal O}\frac1{m^4}
 \right)\\
 &=\frac\pi2+\frac{7\zeta(3)}\pi+\frac{\pi^3}{24}.
\end{aligned}
\]
\end{proof}

It follows from the Rayleigh principle that the optimistic limit of all finite
normalized one-high-leg resonances is at most
\[
        \|M\|.
\]
We next give a fully explicit analytic estimate for this norm.

\subsection{A rigorous bound below \texorpdfstring{$5.466$}{5.466}}

Write $M$ in block form
\[
        M=
        \begin{pmatrix}
        a&b^*\\ b&C
        \end{pmatrix}
\]
relative to $\ell^2(\mathbb N_0)=\R e_0\oplus e_0^\perp$.  Then
\begin{equation}\label{eq:a-c-values}
 a=M_{00}=\frac{14\zeta(3)}\pi,
 \qquad
 c:=\operatorname{tr}C
 =\frac\pi2-\frac{7\zeta(3)}\pi+\frac{\pi^3}{24}.
\end{equation}
Since $C\ge0$,
\[
        \|C\|\le\operatorname{tr}C=c.
\]
Consequently, for $s\in\R$ and $y\in e_0^\perp$,
\[
 \langle M(s,y),(s,y)\rangle
 \le
 a s^2+2\|b\||s|\|y\|+c\|y\|^2.
\]
Hence
\begin{equation}\label{eq:block-upper}
 \|M\|
 \le
 \lambda_{\max}
 \begin{pmatrix}
 a&\|b\|\\ \|b\|&c
 \end{pmatrix}.
\end{equation}

We now bound $\|b\|$ by elementary arithmetic.  For $j\ge1$, partial
fractions give
\begin{equation}\label{eq:b-j-exact}
\begin{aligned}
 b_j=M_{0j}
 &=\frac{16}{\pi}
 \sum_{n\in\mathcal O}\frac1{n(n+2j)^2}\\
 &=\frac1\pi\left(
 \frac{4H_j^{\rm o}}{j^2}
 +\frac{8K_j^{\rm o}}j
 \right)-\frac\pi j,
\end{aligned}
\end{equation}
where
\[
 H_j^{\rm o}=\sum_{r=0}^{j-1}\frac1{2r+1},
 \qquad
 K_j^{\rm o}=\sum_{r=0}^{j-1}\frac1{(2r+1)^2}.
\]
Using the classical bound $333/106<\pi$, the exact rational upper bounds
obtained from \eqref{eq:b-j-exact} for $1\le j\le20$ have squared sum
strictly smaller than $0.574$.  For completeness, upward-rounded values are
listed below:
\[
\begin{array}{c|rrrrrrrrrr}
 j&1&2&3&4&5&6&7&8&9&10\\ \hline
 \overline b_j
 &.678311&.268418&.146878&.093846&.065681
 &.048814&.037858&.030310&.024874&.020818
\end{array}
\]
\[
\begin{array}{c|rrrrrrrrrr}
 j&11&12&13&14&15&16&17&18&19&20\\ \hline
 \overline b_j
 &.017707&.015264&.013308&.011716&.010402
 &.009303&.008375&.007583&.006902&.006311.
\end{array}
\]
The displayed decimals are only for readability; the calculation is rational.
An exact-arithmetic verification script accompanies this note.

For the tail, split the defining sum for $b_j$ at $n=2j$.  This yields
\begin{equation}\label{eq:b-tail-pre}
 b_j
 \le
 \frac4{\pi j^2}\left(H_j^{\rm o}+\frac12\right).
\end{equation}
For $j\ge20$,
\[
        H_j^{\rm o}+\frac12\le\sqrt j.
\]
Indeed, the assertion is a direct rational check at $j=20$; if it holds at
$j$, then
\[
 H_{j+1}^{\rm o}+\frac12
 \le\sqrt j+\frac1{2j+1}
 \le\sqrt{j+1},
\]
because
$\sqrt{j+1}+\sqrt j\le2j+1$.
Thus, for $j\ge21$,
\[
        b_j\le\frac4{\pi j^{3/2}}.
\]
Using again $333/106<\pi$ and the integral test,
\[
\begin{aligned}
 \sum_{j\ge21}b_j^2
 &\le\frac{16}{\pi^2}\sum_{j\ge21}\frac1{j^3}\\
 &<16\left(\frac{106}{333}\right)^2
 \int_{20}^{\infty}x^{-3}\,dx.
\end{aligned}
\]
Combining the finite and tail estimates gives the convenient rational bound
\begin{equation}\label{eq:b-norm-0576}
        \boxed{\|b\|^2<0.576.}
\end{equation}

\begin{theorem}[Rigorous one-high-leg spectral ceiling]
\label{thm:one-leg-5466}
The positive trace-class resonance operator \eqref{eq:resonance-matrix}
satisfies
\begin{equation}\label{eq:5466-ceiling}
        \boxed{\|M\|<5.466.}
\end{equation}
Consequently, every filling-area certificate based on the normalized
one-high-leg hierarchy \eqref{eq:normalized-one-leg} and a global ambient
comass estimate is strictly smaller than $5.466$.
\end{theorem}

\begin{proof}
Besides $333/106<\pi<355/113$, the integral test applied through $n=300$
gives the elementary rational enclosure
\[
        1.202056<\zeta(3)<1.202057.
\]
Hence \eqref{eq:a-c-values} implies
\[
        a<5.356915,
        \qquad
        c<0.184342.
\]
Together with \eqref{eq:b-norm-0576}, it is enough to check that the largest
eigenvalue of
\[
        \begin{pmatrix}
        5.356915&\sqrt{0.576}\\
        \sqrt{0.576}&0.184342
        \end{pmatrix}
\]
is smaller than $5.466$.  This check can be made without evaluating a square
root: with $q=5.466$ one has
\[
        q>5.356915,\qquad q>0.184342,
\]
and
\[
 (q-5.356915)(q-0.184342)>0.576.
\]
Equivalently, $qI$ minus the displayed matrix is positive definite.  The
companion script verifies the slightly sharper rational inequalities before
decimal rounding.
\end{proof}

\begin{remark}[Numerical location of the ceiling]
Finite principal sections suggest that the actual spectral norm is close to
$5.465$.  For example, the largest eigenvalues of sections with a few dozen
resonances stabilize rapidly near this value.  These computations are useful for
orientation, but Theorem~\ref{thm:one-leg-5466} is independent of them.
\end{remark}

\begin{remark}[Scope of the ceiling]
The theorem is deliberately not stated as a ceiling for ``all nonlinear Fourier
methods.''  Arbitrary nonlinear functions of the full Fourier profile contain the
same information as the profile itself and can encode constructions of Z\"ust type.
The obstruction applies to the large but specific class of phase-corrected
\emph{one-high-leg} resonances, and more generally to equivariant corrections for
which a sufficiently high Hardy tangent mode sees only mutually orthogonal shifted
copies of itself.
\end{remark}

\section{Beyond the one-high-leg ceiling}
\label{sec:beyond-ceiling}

The spectral ceiling changes the strategic picture.  Adding higher and higher
powers of $\overline z_1$ cannot close the remaining gap to $2\pi$; one has to use
information which the high-mode test does not see.  We record three concrete
routes.


\subsection{A Fourier--Z\"ust hybrid: current progress}
\label{subsec:hybrid-progress}

Let $\Omega_F$ denote the normalized signed Fourier form and let
$\widetilde\omega_Z$ denote Z\"ust's profile-dependent exact two-form.
Their boundary actions are respectively
\[
        C_F:=\frac{14\zeta(3)}{\pi}
        \qquad\text{and}\qquad
        C_Z:=2\pi.
\]
For $t\in[0,1]$ consider the exact hybrid
\begin{equation}\label{eq:hybrid-form}
        \Theta_t=(1-t)\Omega_F+t\widetilde\omega_Z.
\end{equation}
Its boundary action is
\begin{equation}\label{eq:hybrid-action}
        C(t)=C_F+t(C_Z-C_F).
\end{equation}
Hence any nonzero $t$ for which
\begin{equation}\label{eq:hybrid-target}
        \|\Theta_t\|_{\rm ir}\le1
\end{equation}
holds globally gives a strict improvement over the linear Fourier
constant.  For instance,
\[
        C(1/2)
        =\frac12\left(
        \frac{14\zeta(3)}{\pi}+2\pi
        \right)
        =5.8199788\ldots.
\]
Thus the hybrid has substantially more numerical room than the polynomial
one-high-leg resonance hierarchy.

The calculations below explain why \eqref{eq:hybrid-target} is plausible
locally and isolate the remaining global inequality.  The exact
north-pole diagonalization is rigorous.  The higher-order profile
expansions have been checked algebraically mode-by-mode; turning them into
a complete local theorem in the full $L^\infty$ profile space still
requires the uniform remainder estimates developed in Z\"ust's variation
theory.

\subsubsection{Exact diagonalization at the north-pole profile}

Let
\[
        N(\alpha)\equiv\frac{\pi}{2}.
\]
At this profile Z\"ust's coefficient kernel is
\[
        p_{\alpha,\beta}(N)=1.
\]
Let $\gamma:[0,\pi]\to\C$ be an anti-periodic plane path and expand it on
the doubled circle as
\[
        \gamma(\alpha)
        =\sum_{k\in2\Z+1}c_ke^{ik\alpha}.
\]
With the normalization used throughout this note, the two forms satisfy
the exact spectral identities
\begin{align}
 \Omega_F(\gamma)
 &=
 \sum_{k\in2\Z+1}
 \operatorname{sgn}(k)|c_k|^2,
 \label{eq:hybrid-F-diag}\\
 \widetilde\omega_{Z,N}(\gamma)
 &=
 \sum_{k\in2\Z+1}
 \frac{|c_k|^2}{k}.
 \label{eq:hybrid-Z-diag}
\end{align}
Indeed, for a single odd mode $e^{ik\alpha}$,
\[
 \frac1\pi
 \int_0^\pi\int_\alpha^\pi
 \sin(k(\beta-\alpha))\,d\beta\,d\alpha
 =\frac1k.
\]

Consequently,
\begin{equation}\label{eq:hybrid-diag}
 \Theta_{t,N}(\gamma)
 =
 \sum_{k\in2\Z+1}a_k(t)|c_k|^2,
\end{equation}
where
\[
 a_k(t)=
 \begin{cases}
 1-t+t/k,&k>0,\\[1mm]
 -1+t-t/|k|,&k<0.
 \end{cases}
\]
Thus
\[
        |a_{\pm1}(t)|=1,
\]
whereas for every odd $|k|\ge3$,
\begin{equation}\label{eq:hybrid-spectral-gap}
        |a_k(t)|
        \le1-\frac{2t}{3}.
\end{equation}
In particular, for every $0<t<1$, the huge analytic/anti-analytic
inner-function saturation set of the pure Fourier form collapses at $N$ to
the two first-harmonic planes
\[
        e^{i\phi}e^{i\alpha},
        \qquad
        e^{i\phi}e^{-i\alpha}.
\]
This spectral collapse is the main reason the hybrid is more promising
than another polynomial resonance.

\subsubsection{Quadratic profile variation}

Consider an anti-periodic profile perturbation
\[
        f_\eps(\alpha)
        =\frac{\pi}{2}+\eps q(\alpha),
\]
with
\[
 q(\alpha)
 =
 \sum_{\substack{n\ge1\\ n\ {\rm odd}}}
 \left(a_n\cos n\alpha+b_n\sin n\alpha\right).
\]
At the positive first-harmonic path
\[
        \gamma_0(\alpha)=e^{i\alpha},
\]
the first profile variation of Z\"ust's kernel vanishes.  A direct
second-order expansion gives the diagonal quadratic form
\begin{equation}\label{eq:Z-profile-Hessian}
 \widetilde\omega_{Z,f_\eps}(\gamma_0)
 =
 1+\eps^2
 \sum_{\substack{n\ge1\\n\ {\rm odd}}}
 \lambda_n(a_n^2+b_n^2)
 +O(\eps^4),
\end{equation}
where
\begin{equation}\label{eq:lambda-n}
 \lambda_n
 =
 1
 -
 \sum_{j=1}^{(n+1)/2}\frac1{2j-1}
 -
 \sum_{j=1}^{(n-1)/2}\frac1{2j-1}.
\end{equation}
Equivalently,
\begin{equation}\label{eq:lambda-integral}
 \lambda_n
 =
 \frac12\int_0^\pi
 \frac{\cos t\cos(nt)-\cos^2t}{\sin t}\,dt .
\end{equation}
Hence
\[
        \lambda_1=0,
        \qquad
        \lambda_3=-\frac43,
        \qquad
        \lambda_5=-\frac{28}{15},
\]
and in particular
\begin{equation}\label{eq:profile-gap}
        \lambda_n\le-\frac43
        \qquad(n\ge3,\ n\ {\rm odd}).
\end{equation}
Thus all profile directions transverse to the hemisphere family have a
strictly negative quadratic variation.  The sole zero eigenspace consists
of the first harmonics, precisely the infinitesimal motions inside the
hemisphere family.

\subsubsection{Quadratic path variation}

Now perturb only the maximizing path:
\[
 \gamma_\eps(\alpha)
 =
 e^{i(\alpha+\eps\eta(\alpha))},
 \qquad
 \eta(\alpha)
 =
 \sum_{m\ge1}
 \bigl(A_m\cos2m\alpha+B_m\sin2m\alpha\bigr).
\]
At the north-pole profile one obtains
\begin{align}
 \Omega_F(\gamma_\eps)
 &=
 1-\frac{\eps^2}{2}
 \sum_{m\ge1}(A_m^2+B_m^2)+O(\eps^3),
 \label{eq:F-path-Hessian}\\
 \widetilde\omega_{Z,N}(\gamma_\eps)
 &=
 1-\eps^2
 \sum_{m\ge1}
 \frac{2m^2}{4m^2-1}
 (A_m^2+B_m^2)+O(\eps^3).
 \label{eq:Z-path-Hessian}
\end{align}
Therefore
\begin{equation}\label{eq:hybrid-path-Hessian}
 \Theta_{t,N}(\gamma_\eps)
 =
 1-\eps^2
 \sum_{m\ge1}
 \left[
 \frac{1-t}{2}
 +
 t\frac{2m^2}{4m^2-1}
 \right]
 (A_m^2+B_m^2)
 +O(\eps^3).
\end{equation}
Every path direction has a uniform negative quadratic coefficient.  At
this base profile the profile and path variables decouple to second order,
because the first profile derivative of
$p_{\alpha,\beta}(f)$ vanishes at $N$.

\subsubsection{The first-harmonic flat direction at fourth order}

The only quadratic profile degeneracy is the first-harmonic direction.
Set
\[
        f_r(\alpha)
        =\frac{\pi}{2}+r\cos\alpha,
\]
and allow the maximizing path to move at its natural $r^2$ scale,
\[
        \gamma_{r,b}(\alpha)
        =
        e^{i(\alpha+r^2b\sin2\alpha)}.
\]
Direct fourth-order expansion gives
\begin{equation}\label{eq:Z-quartic}
 \widetilde\omega_{Z,f_r}(\gamma_{r,b})
 =
 1+r^4
 \left(
 -\frac1{24}
 +\frac b3
 -\frac23b^2
 \right)
 +O(r^6),
\end{equation}
while
\begin{equation}\label{eq:F-quartic}
 \Omega_F(\gamma_{r,b})
 =
 1-\frac12b^2r^4+O(r^6).
\end{equation}
For Z\"ust's form alone, the quartic coefficient in
\eqref{eq:Z-quartic} is maximized at $b=1/4$ and equals zero, exactly as
expected from motion along the nearby hemisphere family.  For the hybrid,
optimization over $b$ gives
\begin{equation}\label{eq:hybrid-bstar}
        b_*(t)=\frac{t}{3+t}
\end{equation}
and
\begin{equation}\label{eq:hybrid-quartic}
 \max_b\Theta_t
 =
 1-\frac{t(1-t)}{8(3+t)}\,r^4+O(r^6).
\end{equation}
Thus for every $0<t<1$ the quartic coefficient is strictly negative.
Modulo uniform control of the remainders, the joint quadratic/quartic
calculation strongly indicates that $\Theta_t$ is an actual local
calibration near the hemisphere for every fixed $0<t<1$.

\subsubsection{The Hardy--Z\"ust inequality}

The global problem can be stated in a particularly clean form.  For a
unit-modulus anti-periodic plane path
\[
        \gamma(\alpha)
        =\sum_{k\in2\Z+1}c_ke^{ik\alpha},
\]
the positive-orientation Fourier defect is
\begin{equation}\label{eq:hardy-defect-hybrid}
        1-\Omega_F(\gamma)
        =
        2\sum_{k<0}|c_k|^2
        =
        2\|P_-\gamma\|_{L^2}^2
\end{equation}
under the normalized Fourier convention.  This suggests the following
central inequality.

\begin{question}[Hardy--Z\"ust inequality]\label{q:hardy-zust}
Is it true that for every admissible profile $f\in E(S^1)$ and every
unit-modulus anti-periodic path $\gamma$,
\begin{equation}\label{eq:hardy-zust}
 \boxed{
 \widetilde\omega_{Z,f}(\gamma)-1
 \le
 2\|P_-\gamma\|_{L^2}^2 ?}
\end{equation}
\end{question}

If \eqref{eq:hardy-zust} holds, then
\[
        \frac12(\Omega_F+\widetilde\omega_Z)\le1
\]
pointwise, and the half-hybrid gives immediately
\[
        \Area(M)
        \ge
        \frac12\left(
        \frac{14\zeta(3)}{\pi}+2\pi
        \right)
        =
        5.8199788\ldots.
\]
More generally, it would already be valuable to prove
\begin{equation}\label{eq:hardy-zust-kappa}
 \widetilde\omega_{Z,f}(\gamma)-1
 \le
 \kappa\bigl(1-\Omega_F(\gamma)\bigr).
\end{equation}
Then every
\[
        t\le\frac1{1+\kappa}
\]
is admissible in the hybrid and yields the explicit bound
\[
        \Area(M)
        \ge
        C_F+
        \frac{2\pi-C_F}{1+\kappa}.
\]
Even $\kappa=3$ would already give a bound about $5.5884$, well beyond
the one-high-leg ceiling.

A necessary first subproblem is the positive Hardy restriction
\begin{equation}\label{eq:positive-hardy-target}
 \boxed{
 \widetilde\omega_{Z,f}(\gamma)\le1
 \quad
 \text{for every analytic odd inner path }\gamma.}
\end{equation}
Indeed, $\Omega_F(\gamma)=1$ on these paths.  The local calculations above
show that \eqref{eq:positive-hardy-target} is stable at the hemisphere.

\subsubsection{A nonlocal scalar reduction for the first harmonic}

For the simplest analytic inner path
\[
        \gamma(\alpha)=e^{i\alpha},
\]
one can rearrange Z\"ust's kernel algebraically to obtain
\begin{equation}\label{eq:first-harmonic-nonlocal}
 \widetilde\omega_{Z,f}(e^{i\alpha})
 =
 \frac1\pi\int_0^\pi\csc^2f(\alpha)\,d\alpha
 -
 \frac1{2\pi}
 \int_0^\pi\int_0^\pi
 \frac{
 \bigl(
 \cos f(\alpha)
 -
 \cos(\beta-\alpha)\cos f(\beta)
 \bigr)^2
 }{
 |\sin(\beta-\alpha)|
 \sin^2f(\alpha)\sin^2f(\beta)
 }
 \,d\alpha\,d\beta ,
\end{equation}
with the diagonal interpreted in the limiting sense.
Thus the desired estimate
$\widetilde\omega_{Z,f}(e^{i\alpha})\le1$
is equivalent to a concrete nonlocal Poincar\'e-type inequality.  This
formulation appears more tractable than a pointwise estimate on
$p_{\alpha,\beta}$; in particular, a natural pointwise geometric-mean
bound for the kernel is false.

\subsubsection{Cross-ratio and welding variables}

Assume temporarily that $\operatorname{Lip}(f)<1$ and define
\begin{equation}\label{eq:LR}
        L(\alpha)=\alpha-f(\alpha),
        \qquad
        R(\alpha)=\alpha+f(\alpha).
\end{equation}
Both maps are strictly increasing.  The anti-periodicity
$f(\alpha+\pi)=\pi-f(\alpha)$ implies that
\[
        h:=R\circ L^{-1}
\]
defines an increasing lift of a fixed-point-free circle involution,
satisfying
\begin{equation}\label{eq:h-involution}
        h(h(x))=x+2\pi.
\end{equation}
For $\beta=\alpha+\delta$, introduce
\begin{equation}\label{eq:X-crossratio}
 X_\delta(\alpha)
 =
 \frac{
 \sin\frac{\delta+f(\alpha+\delta)-f(\alpha)}{2}
 \sin\frac{\delta+f(\alpha)-f(\alpha+\delta)}{2}
 }{
 \sin f(\alpha)\sin f(\alpha+\delta)
 }.
\end{equation}
Then a direct trigonometric factorization gives
\begin{equation}\label{eq:p-X}
 \boxed{
 p_{\alpha,\alpha+\delta}(f)\sin\delta
 =
 \frac{4X_\delta(\alpha)
       (1-X_\delta(\alpha))}
      {\sin\delta}.}
\end{equation}
The quantity $X_\delta(\alpha)$ is a circular cross-ratio expression for
the four points
\[
        x,\ y,\ h(x),\ h(y),
        \qquad
        x=L(\alpha),\quad y=L(\alpha+\delta).
\]
Thus the positive-Hardy problem is naturally connected with conformal
welding.

Recent work of Fan--Viklund--Wang introduces an operator from the Fourier
coefficients of the logarithmic chord-distortion kernel
\[
        (z,w)\longmapsto
        \log\left|
        \frac{\varphi(z)-\varphi(w)}{z-w}
        \right|
\]
of a welding homeomorphism and proves a Grunsky-type contraction
inequality.  This makes the following program concrete: determine whether
the Hardy compression of Z\"ust's skew kernel is a positive transform, a
mixed derivative, or a Schur complement of the welding/Grunsky operator
associated with $h$.  Such an identity would potentially imply
\eqref{eq:positive-hardy-target} from a classical or generalized Grunsky
inequality.

At present this welding identification is a research direction, not a
proved equivalence.  The exact relations already established here are
\eqref{eq:LR}--\eqref{eq:p-X} and the de Sitter formulas in
Proposition~\ref{prop:desitter-identities} below.

\subsection{A de Sitter reformulation of the profile geometry}

There is an intrinsic way to see why Z\"ust's kernel contains information which the
polynomial resonances miss.  Let $f:(0,\pi)\to(0,\pi)$ be an anti-periodic
$1$-Lipschitz profile and define
\begin{equation}\label{eq:desitter-map}
 X_f(\alpha)
 =\left(
 \frac{\cos\alpha}{\sin f(\alpha)},
 \frac{\sin\alpha}{\sin f(\alpha)},
 \cot f(\alpha)
 \right)\in\R^{2,1},
\end{equation}
where
\[
        \langle x,y\rangle_{2,1}=x_1y_1+x_2y_2-x_3y_3.
\]
Then
\[
        \langle X_f(\alpha),X_f(\alpha)\rangle_{2,1}=1,
\]
so $X_f$ lies on the two-dimensional de Sitter quadric.

\begin{proposition}[de Sitter identities]\label{prop:desitter-identities}
At almost every $\alpha$,
\begin{equation}\label{eq:desitter-speed}
 \boxed{
 \langle X_f'(\alpha),X_f'(\alpha)\rangle_{2,1}
 =\frac{1-f'(\alpha)^2}{\sin^2f(\alpha)}.}
\end{equation}
For $\alpha\ne\beta\pmod\pi$,
\begin{equation}\label{eq:desitter-inner}
 \langle X_f(\alpha),X_f(\beta)\rangle_{2,1}
 =
 \frac{
 \cos(\beta-\alpha)-\cos f(\alpha)\cos f(\beta)
 }{
 \sin f(\alpha)\sin f(\beta)
 }.
\end{equation}
Moreover, Z\"ust's coefficient kernel satisfies the exact algebraic identity
\begin{equation}\label{eq:zust-desitter-kernel}
 \boxed{
 p_{\alpha,\beta}(f)
 =
 \frac{
 1-\langle X_f(\alpha),X_f(\beta)\rangle_{2,1}^2
 }{
 \sin^2(\beta-\alpha)
 }.}
\end{equation}
\end{proposition}

\begin{proof}
The first two identities follow by direct differentiation and expansion.
For the last one, abbreviate
$d=\beta-\alpha$, $A=\cos f(\alpha)$, and
$B=\cos f(\beta)$.  Then
\[
 1-\langle X_f(\alpha),X_f(\beta)\rangle_{2,1}^2
 =
 \frac{
 1-\cos^2d-A^2-B^2+2AB\cos d
 }{
 \sin^2f(\alpha)\sin^2f(\beta)
 }.
\]
Comparison with the defining formula for $p_{\alpha,\beta}(f)$ gives
\eqref{eq:zust-desitter-kernel}.
\end{proof}

Thus an admissible profile is a spacelike graph on de Sitter space, and the
kernel is a normalized pairwise de Sitter chord distortion.  Hemisphere profiles
are exactly the intersections of the de Sitter quadric with planes through the
origin: indeed, if $\cos h(\alpha)=a_1\cos\alpha+a_2\sin\alpha$, then
\[
        X_h^3=a_1X_h^1+a_2X_h^2.
\]
This reformulates the nonlinear problem as an integral-geometric inequality for
antipodally symmetric spacelike de Sitter curves.  It is qualitatively different
from a polynomial Fourier correction and is a plausible place to seek a much
larger gain.

\subsection{Genuinely multi-high-leg corrections}

The high-mode test behind Theorem~\ref{thm:one-leg-5466} can only be evaded if a
high Hardy tangent mode no longer produces an orthogonal list of shifted copies.
This requires nonlinearities with more than one frequency-dependent ``high leg.''
Such terms are possible, but the boundary coefficients satisfy
$\rho_n\asymp n^{-2}$, so every additional high leg costs two further powers of
$n$.  They therefore have rapidly decaying boundary tails and are unlikely to
recover the missing $0.8$ units by a tail mechanism alone.  Their natural role is
instead a finite-dimensional optimization of the first few modes.  This suggests a
computer-assisted semidefinite search over low-mode equivariant polynomials,
followed by a rigorous interval or sum-of-squares certificate if a worthwhile gain
is found.

The broader lesson is that the remaining gap is likely to require either
profile-adapted geometry, as in \eqref{eq:hybrid-form} and
Proposition~\ref{prop:desitter-identities}, or information coupling the distance
profiles at different interior points (cut loci, geodesic order, or Jacobi-field
constraints).  Purely pointwise one-high-leg Fourier algebra is now quantitatively
exhausted.

\section{Equality structure and Hardy-space defect}

The complex notation also gives an exact interpretation of the loss in the
linear oriented calibration.

Let
\[
        P=\sum_{n>0,\ n\ {\rm odd}}|c_n|^2,\qquad
        N=\sum_{n<0,\ n\ {\rm odd}}|c_n|^2.
\]
If the odd profile derivative has full pointwise norm $|W|=1$, then
$P+N=1$ and
\[
        \Omega_0=P-N.
\]
Consequently
\begin{equation}\label{eq:Hardy-defect}
        \boxed{
        1-\Omega_0=2N}
\end{equation}
for a positive orientation.  Without the equality $|W|=1$, the exact
identity becomes
\[
        1-\Omega_0
        =
        \left(1-\sum_k|c_k|^2\right)
        +2N.
\]
Thus the linear deficit has exactly two components:
\begin{enumerate}[label=(\roman*)]
\item loss of total eikonal energy in the odd projection;
\item negative Hardy-space energy.
\end{enumerate}
The antipodal slack theorem quantifies the first component geometrically,
while the resonant construction attacks the second through autocorrelation
cancellation.

\section{How the method compares with Z\"ust's approach}

The Fourier method and Z\"ust's almost calibration share the same broad
architecture: both encode a filling by functions on the boundary circle and
bound area using a two-form.  The distinction is the level of nonlinearity.

The linear Fourier form is translation invariant in the coefficient space.
Its global constant is explicit and its proof is elementary.  The resonant
form above is still explicit, but its coefficients depend polynomially on
the current profile through $z_1$ and $z_{n+2}$.

Z\"ust constructs a much more geometric point-dependent form directly on
the injective hull $E(S^1)$.  Its comass is exactly one on the hemisphere
and stationary there; this yields a near-sharp second-order estimate for
fillings close to the hemisphere \cite{Zuest2025}.  The 49-page length of
that work reflects this variational stability analysis, not merely the
global numerical bound $\pi^2/2$ recorded separately in its final section.

The present nonlinear resonance can be viewed as a very explicit
finite-order model of the same philosophy:
\[
        \text{adapt the calibration to the profile}
\]
while retaining enough Fourier structure for a global estimate.

\section{What remains open}

The current rigorous constants are
\[
        5.35677234\ldots
        \quad\text{without orientation},
\]
and
\[
        5.38982446\ldots
        \quad\text{with orientation}.
\]
The conjectural value remains
\[
        2\pi=6.28318530\ldots.
\]

Several concrete subproblems emerge.

\begin{question}[Beyond the one-high-leg ceiling]
Theorem~\ref{thm:one-leg-5466} shows that the normalized one-high-leg resonant
hierarchy cannot reach even $5.466$.  Can a Fourier--Z\"ust hybrid
\eqref{eq:hybrid-form}, a de Sitter comparison inequality based on
Proposition~\ref{prop:desitter-identities}, or a genuinely multi-high-leg
low-mode correction produce a substantially larger constant?
\end{question}

\begin{question}[Universal nonlinear improvement]
Can the resonant deformation be combined with the orthogonal untwisting so
that the nonlinear improvement survives without orientability?  There is a
concrete obstruction to simply repeating the oriented proof.  For the
unsigned Jacobian estimate the relevant quadratic energy is
\[
 E_0=\frac12\bigl(\|p\|_2^2+\|q\|_2^2\bigr)
    =\sum_{k\in\mathbb Z}|c_k|^2\le1.
\]
Thus its saturation set consists of all unit-modulus anti-periodic fields
$|W|=1$, not only the analytic and anti-analytic inner fields.  At a mixed
Hardy-spectrum saturator the first variation of the resonant energy contains
terms involving $c_{-1}\overline{c_1}$ and the profile tail
$(z_{n+2})$, and these do not vanish in general.  The cancellation behind
Corollary~\ref{cor:stationary} is therefore genuinely signed.  A universal
nonlinear improvement needs either a different resonance with stationarity
on the full unit-modulus saturation set, or an argument that avoids a
pointwise unsigned-energy relaxation.
\end{question}

\begin{question}[Using the slack defect]
Can the Dirichlet defect in Theorem~\ref{thm:slack} be coupled to a
Z\"ust-type nonlinear form on the anti-periodic projection to obtain a
larger universal constant?
\end{question}

\begin{question}[Global comass of Z\"ust's form]
Z\"ust's numerical computations suggest that his form is not an exact
global calibration but is not far from one.  Any rigorous global comass
bound sufficiently close to one would immediately give a strong global
filling-area estimate.  It would be interesting to combine such a bound
with the Fourier baseline rather than use either form in isolation.
\end{question}

\section{Literature and priority}

Briggs and Wells state that their 2026 orientation-free estimate
$(\sqrt3/4)\pi^2$ appears to be the first quantitative lower bound applying
to arbitrary isometric fillings \cite{BriggsWells2026}.  Z\"ust's 2025
paper explicitly records $\pi^2/2$ as a global lower bound in the oriented
current setting and develops the much deeper almost-calibration theorem
\cite{Zuest2025}.  Targeted searches for the constants
$14\zeta(3)/\pi$ and the Fourier-resonant construction above did not locate
earlier occurrences in the filling-area literature.

Accordingly, the appropriate wording at this stage is that these estimates
are new \emph{to our knowledge}.  A definitive priority determination
should include direct communication with experts working on filling
geometry.

\section{Scaling}

If the boundary circumference is $\ell$, rescaling lengths by
$2\pi/\ell$ gives from Theorem~\ref{thm:universal}
\[
        \Area(M)
        \ge
        \frac{7\zeta(3)}{2\pi^3}\ell^2
\]
without orientability.  The nonlinear oriented estimate scales in the same
quadratic way.

\begin{thebibliography}{99}

\bibitem{BangertCrokeIvanovKatz2005}
V.~Bangert, C.~B.~Croke, S.~V.~Ivanov, and M.~G.~Katz,
\newblock Filling area conjecture and ovalless real hyperelliptic surfaces,
\newblock \emph{Geom. Funct. Anal.} \textbf{15} (2005), no.~3, 577--597.

\bibitem{BriggsWells2026}
J.~Briggs and C.~Wells,
\newblock A discrete view of Gromov's filling area conjecture,
\newblock arXiv:2602.17859, 2026.

\bibitem{Cossarini2020}
M.~Cossarini,
\newblock Discrete surfaces with length and area and minimal fillings of the circle,
\newblock arXiv:2009.02415, 2020.

\bibitem{EvansGariepy}
L.~C.~Evans and R.~F.~Gariepy,
\newblock \emph{Measure Theory and Fine Properties of Functions},
\newblock revised edition, CRC Press, 2015.


\bibitem{FanViklundWang2026}
S.~Fan, F.~Viklund, and Y.~Wang,
\newblock On the Loewner energy of a welding homeomorphism,
\newblock arXiv:2604.16737, 2026.

\bibitem{Gromov1983}
M.~Gromov,
\newblock Filling Riemannian manifolds,
\newblock \emph{J. Differential Geom.} \textbf{18} (1983), no.~1, 1--147.

\bibitem{Hatcher}
A.~Hatcher,
\newblock \emph{Algebraic Topology},
\newblock Cambridge University Press, 2002.

\bibitem{Zuest2025}
R.~Z\"ust,
\newblock The Riemannian hemisphere is almost calibrated in the injective hull of its boundary,
\newblock \emph{Calc. Var. Partial Differential Equations} \textbf{64} (2025),
Art.~297.
\newblock DOI: 10.1007/s00526-025-03133-z.

\end{thebibliography}

\end{document}
